% !TEX root = ../algebra_02.tex

\section{Алгебраические и трансцендентные элементы. Минимальный многочлен алгебраического элемента}

Пусть дано расширение $L/K$ и элемент $a \in L$.

\begin{Definition}{Алгебраический и трансцендентный элементы}{}
	Элемент $a$ называется \textbf{алгебраическим} над $K$, если существует ненулевой многочлен $f \in K[x]$ такой, что $f(a) = 0$.
	В противном случае элемент $a$ называется \textbf{трансцендентным} над $K$.
\end{Definition}

Если элемент $a$ алгебраичен над $K$, рассмотрим множество всех многочленов из $K[x]$, зануляющих $a$:
\[ I = \{ f \in K[x] \mid f(a) = 0 \} \]
Множество $I$ является идеалом в кольце $K[x]$. Так как $K[x]$ --- кольцо главных идеалов, идеал $I$ порождается одним элементом: $I = (f_0)$.

\begin{Definition}{Минимальный многочлен}{}
	Многочлен $f_0$, порождающий идеал $I$, называется \textbf{минимальным многочленом} элемента $a$ над $K$.
	Он определён однозначно с точностью до умножения на обратимый элемент $\varepsilon \in K^*$. Обозначается $\operatorname{Irr}_K(a)$.
\end{Definition}

\begin{Lemma}{Неприводимость минимального многочлена}{}
	Минимальный многочлен $\operatorname{Irr}_K(a)$ неприводим.
\end{Lemma}

\begin{proof}
	Обозначим $f_0 := \operatorname{Irr}_K(a)$. Предположим, что $f_0 = g \cdot h$.
	Так как $f_0(a) = 0$, то $g(a) \cdot h(a) = 0$. Поскольку мы работаем в поле, делителей нуля нет, следовательно, либо $g(a) = 0$, либо $h(a) = 0$.
	Это означает, что либо $f_0 \mid g$, либо $f_0 \mid h$.
	Так как степени $g$ и $h$ не превосходят степени $f_0$, это возможно только в том случае, когда $g \sim f_0$ или $h \sim f_0$.
	Значит, многочлен $f_0$ неприводим.
\end{proof}

\begin{Definition}{Степень алгебраического элемента}{}
	Если элемент $a$ алгебраичен над $K$, то \textbf{степенью} элемента $a$ называется степень его минимального многочлена: $\deg(\operatorname{Irr}_K(a))$.
\end{Definition}

\begin{Proposition}{Степень элемента в конечном расширении}{}
	Пусть $[L:K] = n < \infty$. Тогда любой элемент $a \in L$ является алгебраическим над $K$, и его степень не превосходит $n$.
\end{Proposition}

\begin{proof}
	Рассмотрим набор из $n+1$ элементов: $1, a, a^2, \dots, a^n$.
	Так как размерность пространства $\dim_K L = n$, этот набор линейно зависим. Следовательно, существуют коэффициенты $\alpha_0, \alpha_1, \dots, \alpha_n \in K$, не все равные нулю, такие что:
	\[ \alpha_0 + \alpha_1 a + \dots + \alpha_n a^n = 0 \]
	Рассмотрим многочлен $f(x) := \alpha_n x^n + \dots + \alpha_1 x + \alpha_0$. Так как не все $\alpha_i$ равны нулю, $f \ne 0$. При этом $f(a) = 0$, что означает, что $a$ --- алгебраический элемент над $K$.

	Так как $f(a) = 0$, многочлен $f$ принадлежит идеалу, порожденному $\operatorname{Irr}_K(a)$.
	Следовательно, $\operatorname{Irr}_K(a)$ делит $f$, откуда получаем:
	\[ \deg(\operatorname{Irr}_K(a)) \le \deg(f) \le n \]
\end{proof}
